% \documentclass[10pt, twocolumn, a4paper]{article}
\documentclass[12pt,a4paper]{article}

\usepackage[backend=biber, style=ieee]{biblatex}                        % To include the bibliography
\usepackage[left=2cm, right=2cm, top=2.5cm, bottom=2.5cm]{geometry}     % To set the margins
\usepackage[noend]{algpseudocode}
\usepackage[table]{xcolor}                                              % For coloring cells

\usepackage{algorithm}                                                  % To include algorithms
\usepackage{amsfonts}                                                   % To include math fonts:ToggleTerm direction=float
\usepackage{amsmath}                                                    % To include Mathematic symbols
\usepackage{authblk}                                                    % To format author affiliations
\usepackage{caption}                                                    % For caption spacing
\usepackage{enumitem}                                                   % To customize lists (items like i, ii, iii, iv)
\usepackage{float}                                                      % To place figures where you want them
\usepackage{fancyhdr}                                                   % To customize headers and footers
\usepackage{graphicx}                                                   % To include images
\usepackage{hyperref}                                                   % To include hyperlinks
\usepackage{lipsum}                                                     % TODO: remove this
\usepackage{listings}                                                   % To include code
\usepackage{tabularx}                                                   % For equal-width columns
\usepackage{tcolorbox}                                                  % To make colored boxes
\usepackage{tikz}                                                       % To draw graphs
\usepackage{titlesec}                                                   % To format section titles
\usepackage{xcolor}                                                     % To define colors

\usetikzlibrary{graphs,graphs.standard}
\usetikzlibrary{positioning}

\addbibresource{./references.bib}

% Esto es para poder hacer cajitas de código con el fondo gris
\lstset{
    language=C++,
    basicstyle=\ttfamily\footnotesize,
    keywordstyle=\color{blue}\bfseries,
    stringstyle=\color{green!60!black},
    commentstyle=\color{gray},
    backgroundcolor=\color{gray!05},
    frame=single,
    numbers=left,
    numberstyle=\footnotesize,
    stepnumber=1,
    numbersep=10pt,
    tabsize=2,
    showstringspaces=false,
    captionpos=b,
    breaklines=true,
}

% Para poder hacer flechas
\usetikzlibrary{shapes, arrows}

% Sección de definiciones
\titleformat{\section}{\Large\bfseries}{\thesection}{1em}{}
\titleformat{\subsection}{\large\bfseries}{\thesubsection}{1em}{}

% Caja de colores
\definecolor{mint}{RGB}{202,251,202}
\definecolor{yellow}{RGB}{255,255,202}
\definecolor{red}{RGB}{255,202,202}

% Variables globales para el documento
\newcommand{\facultyname}{UP | Facultad de Ingeniería}
\newcommand{\coursename}{Programación con GPUs}
\newcommand{\currentsemester}{Segundo Semestre}
\newcommand{\currentyear}{2025}

% Esto es para poder agregar comentarios al código
\newcommand{\comentario}[1]{\textcolor{gray}{// #1}}

% Definir los encabezados y pies de página
\pagestyle{fancy}
\fancyhf{} % Borra encabezados y pies de página

% Header configuration
\fancyhead[R]{\includegraphics[height=50px]{../apuntes/images/logo_up.jpg}}
\renewcommand{\headrulewidth}{0.4pt} % Agregar línea debajo del encabezado
\setlength{\headheight}{50pt} % Ajustar la altura del encabezado

% Footer configuration
\fancyfoot[L]{\facultyname}
\fancyfoot[C]{\thepage}
\fancyfoot[R]{\coursename}
\renewcommand{\footrulewidth}{0.4pt} % Agregar línea arriba del pie de página

% Configuración de la primera página
\fancypagestyle{firstpage}{
  \fancyhf{} % Borra encabezados y pies de página
  \fancyfoot[C]{\thepage} % Número de página centrado
  \renewcommand{\headrulewidth}{0pt} % Sin línea en el encabezado
  \renewcommand{\footrulewidth}{0.4pt} % Sin línea en el pie de página
}

% Configuración de la primera página
\AtBeginDocument{\thispagestyle{firstpage}}


\begin{document}

\begin{center}
  \LARGE\textbf{\coursename} \\
  \Large{Práctica 02 - Introducción a CUDA} \\
  \normalsize{\currentsemester, \currentyear} \\
  \vspace{1em}
  \hrule
\end{center}

\setcounter{section}{2}

\input{../shared.tex/practica_warning_box.tex}

\newpage

\tableofcontents

\newpage

\subsection{Ejercicio: Suma de dos vectores}

Si quisiéramos utilizar un thread para calcular la suma de los dos vectores. ¿Cómo se calcularía el índice del thread?

\begin{enumerate}[label=\roman*.]
  \item \texttt{idx = threadIdx.x + threadIdx.y;}
  \item \texttt{idx = blockIdx.x + threadIdx.x;}
  \item \texttt{idx = blockIdx.x * blockDim.x + threadIdx.x;}
  \item \texttt{idx = blockIdx.x * threadIdx.x;}
\end{enumerate}

\subsection{Ejercicio: Sumar elementos adyacentes}

Supongamos ahora que queremos usar un thread para sumar dos elementos contiguos de un vector (por simplicidad podemos
suponer que el vector tiene un número par de elementos). ¿Cómo se calcularía el índice del thread?

\begin{enumerate}[label=\roman*.]
  \item \texttt{idx = blockIdx.x * blockDim.x + threadIdx.x + 2;}
  \item \texttt{idx = blockIdx.x * threadIdx.x * 2;}
  \item \texttt{idx = (blockIdx.x * blockDim.x + threadIdx.x) * 2;}
  \item \texttt{idx = blockIdx.x * blockDim.x * 2 + threadIdx.x;}
\end{enumerate}

\subsection{Ejercicio: Calcular el tamaño de los threads}

Para un programa que suma dos vectores, supongamos que el tamaño de los vectores es 8000 y que el tamaño del bloque de
threads es 1024. El programador configura que el kernel lance la menor cantidad de bloques de threads para cubrir todos
los elementos de salida. ¿Cuántos threads habrá corriendo en el \textit{grid}?

\subsection{Ejercicio: cudaMalloc (parte 1)}

Si queremos reservar un array de \texttt{v} elementos de tipo \texttt{int} en la memoria global del dispositivo CUDA,
¿cuál sería la expresión apropiada para el segundo argumento de la llamada a \texttt{cudaMalloc}?

\subsection{Ejercicio: cudaMalloc (parte 2)}

Si queremos reservar un array de \texttt{n} elementos del tipo \texttt{float} y tenemos una variable \texttt{d\_A} que
apunta a la memoria reservada, ¿cuál sería la expresión apropiada para el primer argumento de la llamada a
\texttt{cudaMalloc}?

% cudaMemcpy
\subsection{Ejercicio: Copia de memoria desde el \textit{host} al \textit{device}}

Si queremos copiar 3000 bytes de datos desde el array del host \texttt{h\_A} (donde \texttt{h\_A} es un puntero al
elemento 0 del array de origen) al array del dispositivo \texttt{d\_A} (donde \texttt{d\_A} es un puntero al elemento 0
del array de destino), ¿cuál sería la llamada a la API apropiada para esta copia de datos en CUDA?

% Error handling
\subsection{Ejercicio: Manejo de errores}

¿Cómo declararías una variable \texttt{err} que pueda recibir apropiadamente el valor devuelto por una llamada a la API
de CUDA en caso de error?

\subsection{Ejercicio: funciones en CUDA}

Estás trabajando en un proyecto de CUDA y uno de tus compañeros se uqeja de que CUDA es muy tedioso ya que hubo que
declarar muchas funciones que se planean ejecutar tanto en el host como en el device dos veces, una como función de host
y otra como función de device. ¿Cuál sería tu respuesta?

\input{../shared.tex/common_footers.tex}

\end{document}
